% ===================================================================
%  CADRO N.º 12 — A estación. — Os ferrocarrís. — Os barcos.
%  Traduction 2026 d'apres texto/io, controlee sur texto/fr, texto/es
%  et texto/pt. Consigne : voir texto/gl/10-cadro-01.tex.
%
%  MEME OBSERVATION QU'AU TABLEAU 11, ET LE RAIL LA REND PLUS NETTE
%  ENCORE QUE LA VILLE : « vagón », « anden », « semáforo »,
%  « ténder », « locomotora », « furgón », « telégrafo », « billete »
%  sont venus tout faits avec le chemin de fer, et les trois langues
%  les ont pris ensemble.
%
%  CE QUI RESTE GALICIEN EST CE QUE LE RAIL A TROUVE EN PLACE : les
%  mots de la main et du bagage — « maleta », « caixa », « cesto »,
%  « corda », « roda », « asubío » le sifflet, « faiado » la soute.
%  Le rail n'a pas remplace la valise ; il l'a chargee dans un wagon
%  dont il apportait le nom.
%
%  ZERO INVERSION ATTENDUE, ZERO RELEVEE. Ce tableau enumere plus
%  qu'il ne construit, comme dans les cinq colonnes precedentes, et
%  une liste ne s'inverse pas.
% ===================================================================

%%K t12-apar-1 apar
\VUsaut{4.0mm}
\VUtitre{15.0pt}{60}{CADRO N.º 12}
\VUsaut{2.4mm}
\VUpk{11.4pt}{40}{A estación. --- Os ferrocarrís. --- Os barcos.}
\VUsaut{3.4mm}
\textsc{Nota.} --- \textit{Os horarios empregados en cada estación son distintos; por
iso tomamos como exemplo as horas do} \VUgras{cadro francés}.

%%K t12-01-1 p
1. --- Acabo de viaxar por Alemaña. Antes de saír merquei unha \VUgras{guía de
horarios} \textsuperscript{(15)} para saber a hora de saída dos trens. Despois de
comprobar que o tren máis cómodo era o que saía de París ás nove da noite e chegaba a
Estrasburgo á unha e trece minutos, preparei a miña viaxe e, ao día seguinte,
fíxenme levar á estación.

%%K t12-02-1 p
2. --- Cando entrei no gran vestíbulo de saídas, detrás dun vello \VUgras{crego}
\textsuperscript{(2)} de aldea que levaba na man o seu \VUgras{saco de viaxe}
\textsuperscript{(3)}, moitos \VUgras{viaxeiros} \textsuperscript{(1)} chegaran xa.

%%K t12-02-2 p
Como quería mandar un \VUgras{despacho (telegrama)} \textsuperscript{(5)}, fixen que
un \VUgras{revisor} \textsuperscript{(6)} me indicase \VUgras{a oficina do telégrafo}
\textsuperscript{(4)}.

%%K t12-03-1 p
3. --- Antes de pasar á \VUgras{sala de espera} \textsuperscript{(7)} fun proverme
dun \VUgras{billete} \textsuperscript{(9)} de ida e volta.

%%K t12-04-1 p
4. --- Non agardei moito na \VUgras{ventaíña} \textsuperscript{(8)}, pois había pouca
xente. Un \VUgras{soldado} \textsuperscript{(10)}, que saía de permiso, tivo a
amabilidade de cederme a súa quenda, de xeito que tiven tempo dabondo para pedir algúns
informes ao \VUgras{xefe de estación} \textsuperscript{(13)}, que estaba a conversar
co \VUgras{comisario} \textsuperscript{(12)} de vixilancia e co \VUgras{garda}
\textsuperscript{(11)} de servizo.

%%K t12-tit-1 sub
\VUpk{10.2pt}{40}{Facturación da equipaxe.}
\VUsaut{3.0mm}

%%K t12-05-1 p
5. --- Despois de mercar un xornal na \VUgras{biblioteca} \textsuperscript{(14)},
fixen pesar a miña \VUgras{equipaxe} \textsuperscript{(17)}, que un \VUgras{mozo}
\textsuperscript{(16)} acababa de traer nun \VUgras{carro de paquetes}
\textsuperscript{(19)}; \VUgras{o empregado} \textsuperscript{(22)} encargado da
\VUgras{báscula} \textsuperscript{(21)} fíxome saber que o meu baúl pesaba trinta e
cinco quilogramos; iso daba un \VUgras{exceso de peso} de cinco quilogramos, polo que
debía pagar un suplemento na \VUgras{ventaíña} da equipaxe \textsuperscript{(23)}.

%%K t12-06-1 p
6. --- Como ía moi cedo, tiña tempo libre para examinar o movemento dos viaxeiros na
estación. Uns deixaban ou recollían a súa equipaxe lixeira na \VUgras{consigna}
\textsuperscript{(25)}; outros consultaban os \VUgras{horarios}
\textsuperscript{(26)}. Os viaxeiros que chegaban apuraban cara á \VUgras{saída}
\textsuperscript{(32)} coa súa \VUgras{maleta} \textsuperscript{(24)} na man. Algúns
agardaban na \VUgras{entrega de equipaxe} \textsuperscript{(27)} e presentaban o seu
\VUgras{billete} \textsuperscript{(29)} para recuperaren os seus baúles, as súas
\VUgras{caixas} \textsuperscript{(28)} ou os seus \VUgras{cestos}
\textsuperscript{(30)}.

%%K t12-07-1 p
7. --- \VUgras{Os empregados de consumos} \textsuperscript{(33)} examinaban con
coidado os paquetes para comprobaren se había algo que declarar.

%%K t12-08-1 p
8. --- Algúns viaxeiros foron á \VUgras{cantina} \textsuperscript{(34)}, onde un
\VUgras{camareiro} \textsuperscript{(35)} se afanaba por servilos. Teñen que apurar,
pois o \VUgras{tren} \textsuperscript{(36)}, que é expreso, non ficará moito na
estación.

%%K t12-09-1 p
9. --- Despois fun ao anden de saída e busquei un \VUgras{vagón}
\textsuperscript{(37)} directo para non me ver obrigado a mudar de vagón durante a
viaxe. Subín a un coche de corredor que ten un \VUgras{departamento}
\textsuperscript{(38)} para fumadores e un departamento reservado a «señoras soas».

%%K t12-tit-2 sub
\VUpk{10.2pt}{40}{Saída de noite.}
\VUsaut{3.0mm}

%%K t12-10-1 p
10. --- Despois de subir polo \VUgras{estribo} \textsuperscript{(40)}, pechar a
\VUgras{porta} \textsuperscript{(39)}, enrolar a \VUgras{cortina}
\textsuperscript{(41)} e colocar a miña equipaxe pequena na \VUgras{rede}
\textsuperscript{(43)}, sentei no \VUgras{banco} \textsuperscript{(42)}. Ao ver o
\VUgras{lampareiro} \textsuperscript{(44)} acender as lámpadas, pensei que tería que
pasar unha noite no vagón e chamei o empregado encargado do alugueiro das
\VUgras{almofadas} \textsuperscript{(45)} para lle pedir unha.

%%K t12-11-1 p
11. --- O revisor do tren veu comprobar os billetes. Pregunteille por que non
saíramos aínda, pois era a hora xusta. Respondeume que o \VUgras{xefe do tren}
\textsuperscript{(46)} estaba ocupado en meter bicicletas no \VUgras{furgón}
\textsuperscript{(47)}. Ademais, aínda non se cambiaran todos os quentapés
\textsuperscript{(48)}.

%%K t12-12-1 p
12. --- Pero axiña iamos saír, pois o \VUgras{fogueiro} \textsuperscript{(50)}, no seu
\VUgras{ténder} \textsuperscript{(49)}, collía carbón para atizar o lume da
\VUgras{locomotora} \textsuperscript{(54)}, e \VUgras{o maquinista}
\textsuperscript{(51)} só agardaba polo sinal do \VUgras{subxefe de estación}
\textsuperscript{(52)}, que estaba no \VUgras{anden} \textsuperscript{(53)}.

%%K t12-13-1 p
13. --- A \VUgras{caldeira} \textsuperscript{(56)} da locomotora roñaba xordamente; a
\VUgras{cheminea} \textsuperscript{(55)} vomitaba torrentes de fume mesturado con
\VUgras{vapor} \textsuperscript{(60)}. Un \VUgras{asubío} estridente
\textsuperscript{(59)} soou con forza e os \VUgras{émbolos} \textsuperscript{(58)}
comezaron a moverse, empurrando as \VUgras{rodas} \textsuperscript{(57)} da pesada
máquina.

%%K t12-tit-3 sub
\VUsaut{3.0mm}
\VUpk{10.2pt}{40}{Saída!}
\VUsaut{3.0mm}

%%K t12-14-1 p
14. --- A velocidade do tren, que corría sobre os \VUgras{carrís}
\textsuperscript{(64)}, foi acelerándose; os \VUgras{andéns} \textsuperscript{(65)}
desapareceron; pasamos diante dos \VUgras{retretes} \textsuperscript{(66)} e tamén
diante dun conxunto de vagóns estacionados na outra \VUgras{vía}
\textsuperscript{(63)}: \VUgras{coches cama} \textsuperscript{(67)}, \VUgras{coche
restaurante} \textsuperscript{(68)}, \VUgras{coche correo} \textsuperscript{(69)}.

%%K t12-noto-1 noto
\VUnotes{143.2mm}{%
\textit{(*) Nota.} --- \textit{Compréndese que a descrición da viaxe se fai segundo a
fronteira de antes da guerra.}%
}

%%K t12-15-1 p
15. --- Despois pasou diante dos nosos ollos a \VUgras{estación de mercadorías}
\textsuperscript{(71)}. Baixo os alpendres, uns empregados cargaban as
\VUgras{mercadorías} \textsuperscript{(72)} expedidas por tren lento. Uns
\VUgras{cadrilleiros} \textsuperscript{(76)} preto do \VUgras{depósito de auga}
\textsuperscript{(73)} manobraban \VUgras{vagóns de gando} \textsuperscript{(75)} e
\VUgras{vagóns plataforma} \textsuperscript{(79)} para formaren un \VUgras{tren de
mercadorías} \textsuperscript{(80)}.

%%K t12-16-1 p
16. --- Pasamos o último \VUgras{cambio de vía} \textsuperscript{(78)}, preto do cal
estaba o guardaagullas; deixamos atrás o \VUgras{disco de sinais}
\textsuperscript{(86)} e a estación desapareceu ao lonxe.

%%K t12-17-1 p
17. --- Axiña cruzamos unha \VUgras{ponte de ferro} \textsuperscript{(74)} que pasa
por riba do \VUgras{río} \textsuperscript{(81)}. Despois de cruzar esa ponte fomos ao
longo da ribeira, na que uns \VUgras{remolcadores} \textsuperscript{(82)} potentes
tiraban de pesadas \VUgras{gabarras} \textsuperscript{(83)}. Vimos ademais un
\VUgras{barco de viaxeiros} \textsuperscript{(84)} cando atracaba nun \VUgras{pontón}
\textsuperscript{(85)}.

%%K t12-18-1 p
18. --- Despois de pasar a toda présa diante de varias estacións, cruzamos algúns
\VUgras{túneles} \textsuperscript{(87)} e deixamos atrás \VUgras{casiñas} sen conta
\textsuperscript{(88)} de \VUgras{gardabarreiras.} Arrolado polo abalo do tren,
quedei profundamente durmido.

%%K t12-tit-4 sub
\VUsaut{3.0mm}
\VUpk{10.2pt}{40}{Estación de Deutsch-Avricourt.}
\VUsaut{3.0mm}

%%K t12-19-1 p
19. --- Espertoume de súpeto a voz dun empregado que berraba: «Deutsch-Avricourt!
\textsuperscript{(*)}. Baixen todos!» Fíxose a inspección da aduana e todas as
formalidades aburridas da fronteira. Tiven que pór o meu reloxo de peto en
concordancia co \VUgras{reloxo grande} \textsuperscript{(89)} da estación, que ía
cincuenta e cinco minutos adiantado; apenas esperto, distinguín a custo a
\VUgras{agulla grande} \textsuperscript{(90)} da \VUgras{pequena}
\textsuperscript{(91)}; a propia \VUgras{esfera} \textsuperscript{(92)} estaba de todo
confusa por mor da néboa da mañá.

%%K t12-20-1 p
20. --- Mentres os aduaneiros facían a súa inspección, eu, a bocexar, descifraba os
\VUgras{carteis} \textsuperscript{(93)} que adornan as paredes da estación. Almorcei
para pasar o tempo, consultei o mapa da \VUgras{rede alemá} \textsuperscript{(94)}
para ver canta distancia me separaba aínda de Estrasburgo, e volvín entrar no meu
vagón, reconfortado pola esperanza de acadar axiña a meta da miña viaxe.
